\documentclass[12pt]{article}
\usepackage[utf8]{inputenc}
\usepackage[T1]{fontenc}
\usepackage{hyperref}
\title{Zambian Sign Language Language Resources for K-12 Curriculum Development}
\author{Compiled by Claude Sonnet 5.5 (claude-sonnet-5-5) for Andr\'as Kornai}
\date{}
\begin{document}
\maketitle

\section{Introduction}
Zambian Sign Language (ISO 639-3 zsl) is the sign language of the deaf community in Zambia. The search results show some dictionary work and several studies of deaf education. They also point to a continuing shortage of standardised vocabulary, teaching materials and trained teachers.

\section{Primary Educational Resources}
\begin{itemize}
\item A. L. Bwalya, \emph{Zambian Sign Language} (1985), from the Curriculum Development Centre of the Ministry of General Education and Culture, according to a Glottolog bibliographic record. I did not see the book itself. \url{https://glottolog.org/resource/reference/id/319313}
\item University of Zambia thesis on Magwero and St.\ Joseph schools for the deaf (Mulonda, 2013). It found that most teachers and pupils want sign language taught as an examinable subject. \url{https://dspace.unza.zm/handle/123456789/7167}
\end{itemize}

\section{Linguistic and Reference Materials}
\begin{itemize}
\item Vincent Chanda, ``Lexicography and Sign Language Engineering: The Zambian Experience'', \emph{Lexikos} (1997). It describes the Zambian Sign Language Dictionary Project. \url{https://lexikos.journals.ac.za/pub/article/view/980}
\item Zambia Association of the Deaf, \emph{Zambian Sign Language Dictionary} (Lusaka, 2001), per a Glottolog record. \url{https://glottolog.org/resource/reference/id/552544}
\item The same Chanda article on the African Journals Online site. \url{https://www.ajol.info/index.php/lex/article/view/146269}
\end{itemize}

\section{Religious and Community Resources}
The Zambia Association of the Deaf appears as publisher of the 2001 dictionary. I found no other community or religious resources in the search results.

\section{Digital and Audio Resources}
No video corpus, online dictionary or signed teaching videos were found in the search results.

\section{Curriculum Development Context}
Studies from the University of Zambia report gaps in sign vocabulary for school subjects, little teacher training in sign language, and few adapted materials. They also report a lack of agreed local signs across the country. The studies include Muzata and Mahlo (2019) on the 2013 revised curriculum, and Zulu (2019) on Lusaka primary schools. A curriculum project would need to coordinate with the Zambia Association of the Deaf and the Ministry of General Education on standard signs. \url{https://dspace.unza.zm/handle/123456789/6485}

\end{document}